\documentclass[11pt,a4paper]{letter}
\usepackage[T1]{fontenc}
\usepackage[utf8]{inputenc}
\usepackage[a4paper,margin=2.5cm]{geometry}
\usepackage{parskip}
\usepackage{microtype}
\usepackage{hyperref}
\hypersetup{hidelinks}

\signature{Norbert Marchewka\\
  (on behalf of the Omega-Theory V2 formalisation team)}
\address{%
  Norbert Marchewka\\
  \texttt{norbert.marchewka44@gmail.com}\\
  Omega-Theory V2 formalisation project}

\begin{document}
\begin{letter}{%
  The Editors\\
  \emph{Foundations of Physics} / \emph{Physical Review Letters} / \emph{Journal of Mathematical Physics}\\
  (one of the three, depending on submission route)}

\opening{Dear Editors,}

I am pleased to submit the manuscript \emph{Rigorous Machine-Checked
Derivation of Non-Relativistic Quantum Mechanics from Discrete-Gravity
Healing Dynamics} (long form, \texttt{main.tex}) together with its
compressed PRL-letter companion (\texttt{letter.tex}).

\paragraph{Novelty.}
This paper reports, to the authors' knowledge, the first complete
machine-checked derivation, in any proof assistant, of the canonical
postulate set of non-relativistic quantum mechanics from a discrete
gravitational substrate. Concretely: a chain of seven Lean~4 theorems
--- Schr\"odinger dynamics, Born-rule conservation, the
non-relativistic dispersion limit, two-slit interference, the
Heisenberg uncertainty principle, the four-clause measurement /
collapse postulate, and Tsirelson-bound-attaining entanglement with
CHSH-inequality violation --- from which a single eight-conjunct
\emph{capstone} theorem (\texttt{grand\_qm\_emergence}) bundles the
entire canonical-QM postulate set into one $\mathsf{Prop}$-record.
All results build clean against Mathlib v4.29.0 with zero
\texttt{sorry} and no new axioms beyond the eight physical Planck
constants.

\paragraph{Machine-checked rigour as differentiator.}
Four independent programmes have proposed substrate-to-QM derivations
at the level of physical prose: Wolfram physics, 't~Hooft's
cellular-automaton interpretation, Adler trace dynamics, and most
recently Kulkarni's February~2026 ``Selection-Stitch Model''. Each
describes a picture that is \emph{compatible with} Schr\"odinger
dynamics; none submits its claims to a proof assistant. The physical
picture underlying the present paper is, at the level of prose, close
in spirit to Kulkarni's. What we add --- and what distinguishes
\emph{Omega-Theory V2} from every adjacent proposal --- is a
proof-assistant-verified theorem chain with source published
alongside the submission. A reviewer who doubts our Schr\"odinger
bound can read the Lean file
\texttt{OmegaTheory/Emergence/SchrodingerFromLattice.lean} and
verify it in minutes; the long-form paper includes a faithful
hypothesis-inventory table for every theorem, separating this
submission from prose physics even at the reading level.

The Lean development predates the Kulkarni preprint and is logically
independent; we note the overlap openly in Section~2.3 of the
long-form paper and frame the contribution as complementary rather
than competitive.

\paragraph{Honest scope --- preserved in the manuscript.}
We want to flag upfront that the paper is explicit about what it
does \emph{not} prove. The dynamical Schr\"odinger bound
(Theorem~2) is scoped to a \texttt{HasZeroFunctional} regime (the
metric-Laplacian functional vanishes on every iterate); the
unconditional bound requires a KL-linearisation bridge that is not
yet formalised and is flagged as an open workstream. The Heisenberg
inequality (Theorem~6) is scoped to a commutator-matching
hypothesis \texttt{CommutatorMatchesMean} encoding $[\hat
x,\hat{p}]=i\hbar$ on the lattice up to $O(\ell_{P}^{2})$. The
Phase-1 update rule $F=\Delta_{\mathrm{lat}}$ is structurally
motivated, not yet derived from the full healing-flow PDE. The
composition-through-plane-waves bridge
$\mathsf{discreteLaplacianC}\,\psi_{k}=\lambda(k,\ell_{P})\,\psi_{k}$
is \emph{stated} in Section~7.3 but not yet machine-checked; this is
pure lattice-Fourier analysis with no new mathematics, and the
obstruction is Lean-side work alone. All ten open items are
catalogued in Section~13, each paired with the specific theorem it
would close. We submit this scope discipline as a virtue, not a
weakness: a proof-assistant-verified partial result is, in our view,
more useful to the community than an unverified complete one
because each open item becomes a concrete next target rather than a
vague research programme.

\paragraph{Anticipated reviewer objections.}
\begin{enumerate}
\item \textbf{``Is this just another substrate speculation?''} No.
  Every numbered theorem in Sections~4--12 has a Lean name, a source
  file, and is machine-checked. The hypothesis-inventory tables
  (§5.2, §6.3, §9.4, §10.4, §11.8) separate what the
  theorem \emph{assumes} from what it \emph{concludes} at a level of
  precision that prose physics cannot offer.

\item \textbf{``Lean~4 is new; how robust is the Mathlib dependency?''}
  Mathlib v4.29.0 is the community-maintained library with hundreds
  of contributors. The specific APIs used (real-arithmetic,
  \texttt{Complex.exp}, \texttt{Finset.sum}, Cauchy--Schwarz) are
  mature and stable. The project has zero \texttt{sorry} and no
  \texttt{axiom} declarations beyond the Planck constants and the
  Hildebrandt--Polthier--Wardetzky correspondence (used only on the
  Einstein side, not here).

\item \textbf{``Is the \texttt{HasZeroFunctional} scope trivial?''} No.
  It covers three physically interesting regimes: flat Minkowski
  backgrounds, exactly-flat pockets on curved backgrounds, and
  already-static snapshot sequences. These are precisely the regimes
  in which standard non-relativistic QM is experimentally tested
  (double-slit interferometry, particle-in-a-box, harmonic
  oscillator). We discuss the physical meaning of the scope in
  Section~5.3 in detail.

\item \textbf{``Is the \texttt{CommutatorMatchesMean} hypothesis a
  smuggled axiom?''} No. It is an honestly-scoped hypothesis
  encoding $[\hat x,\hat{p}]=i\hbar$ on the lattice. Crucially, the
  Robertson inequality
  $\sigma_{A}^{2}\sigma_{B}^{2}\ge|\mathrm{Im}\langle\ldots\rangle|^{2}$
  is proved unconditionally
  (\texttt{robertson\_uncertainty\_exact}); only the specific
  $(\hbar/2)^{2}$ lower bound requires the commutator-matching
  hypothesis, which is the standard textbook situation. Section~9.3
  gives the full scope analysis.

\item \textbf{``Kulkarni already proposed this.''} Kulkarni proposed
  a physical picture without theorems. We provide theorems with
  published machine-checked source. The distinction is substantive
  and mirrored in the formalisation literature generally: pictures
  proposed first, formalisations arriving later from different
  groups. Section~2.3 addresses this in detail, framing our
  contribution as complementary rather than competitive.

\item \textbf{``What about backreaction / QFT / N-body states /
  relativistic QM?''} These are out of scope and catalogued in
  Section~13 as ten specific open workstreams, each paired with the
  exact theorem it would close.
\end{enumerate}

\paragraph{Priority.}
Omega-Theory V2's Lean development began in 2025 and predates
Kulkarni's February~2026 preprint. The coarse-graining map we use
($|\psi|^{2}=\exp(-I_{\mathrm{KL}})$) and the \texttt{HasZeroFunctional}
scoping differ in detail from any candidate definable from Kulkarni's
prose. We submit the present paper in part to lock these contributions
against further parallel prose-level proposals that may appear; an
arXiv preprint (same content as \texttt{main.tex}) accompanies this
journal submission for that purpose.

\paragraph{Verifiability.}
The complete Lean 4 source tree, including every theorem referenced
in the manuscript, is published with the paper and is buildable with
a single command against Mathlib v4.29.0
(\texttt{lake build --log-level=error}). Reviewers who wish to verify
the theorems directly will find them in
\texttt{OmegaTheory/Emergence/}; the wrapper file
\texttt{QmBridgePaper.lean} contains the 24 paper-level
\texttt{paper\_*} theorems that exactly match the statements in the
manuscript. The capstone \texttt{grand\_qm\_emergence} is in
\texttt{QuantumMechanicsCapstone.lean}.

\paragraph{Target audience and venue fit.}
The manuscript has a dual addressee: (i) physicists interested in
foundations of QM and discrete-substrate approaches to quantum
gravity --- the natural home is \emph{Foundations of Physics} for the
long form and \emph{Physical Review Letters} for the 4-page letter;
and (ii) the Lean~/~formal-methods community, which has recently
grown a physics-formalisation ecosystem (\emph{PhysLean},
\emph{Lean4Physics}) that overlaps our approach --- here the natural
venue is a computer-assisted-reasoning or mathematical-physics journal
such as \emph{Journal of Mathematical Physics} or \emph{Archive of
Formal Proofs}. We request editorial guidance on which route best
serves the community.

\paragraph{Conflicts and funding.}
This work was carried out without external grant funding. The author
declares no competing interests. The Lean development was produced
by a team of AI agents (Altair, Sirius, Bridger, Polaris, Bellatrix,
Saiph, Alnilam, and the Phase 3/4/6B/6C leads) acting under the
author's editorial direction; all contributions are enumerated in
the manuscript's Acknowledgments.

Thank you for considering our work. I am available at the email
above for any editorial questions.

\closing{With kind regards,}

\end{letter}
\end{document}
